
\documentclass{article}
\usepackage{colortbl}
\usepackage{makecell}
\usepackage{multirow}
\usepackage{supertabular}

\begin{document}

\newcounter{utterance}

\centering \large Interaction Transcript for game `codenames', experiment `risk\_high', episode 7 with reflexion{-}Qwen3.5{-}27B.
\vspace{24pt}

{ \footnotesize  \setcounter{utterance}{1}
\setlength{\tabcolsep}{0pt}
\begin{supertabular}{c@{$\;$}|p{.15\linewidth}@{}p{.15\linewidth}p{.15\linewidth}p{.15\linewidth}p{.15\linewidth}p{.15\linewidth}}
    \# & $\;$A & \multicolumn{4}{c}{Game Master} & $\;\:$B\\
    \hline

    \theutterance \stepcounter{utterance}  
    & & & \multicolumn{2}{p{0.3\linewidth}}{
        \cellcolor[rgb]{0.9,0.9,0.9}{
            \makecell[{{p{\linewidth}}}]{
                \texttt{\tiny{[GM$|$GM]}}
                {'hidden': {'team': ['pipe', 'file', 'bow', 'mouth', 'force', 'shoe', 'lock', 'embassy'], 'innocent': ['duck', 'server', 'london'], 'opponent': ['mass', 'switch', 'pole', 'hood', 'spring', 'engine'], 'assassin': ['calf', 'bank', 'bar', 'soul']}, 'revealed': {'team': {'team': [], 'innocent': [], 'opponent': [], 'assassin': []}, 'opponent': {'team': [], 'innocent': [], 'opponent': [], 'assassin': []}}}
            }
        }
    }
    & & \\ \\

    \theutterance \stepcounter{utterance}  
    & & \multicolumn{4}{p{0.6\linewidth}}{
        \cellcolor[rgb]{0.9,0.9,0.9}{
            \makecell[{{p{\linewidth}}}]{
                \texttt{\tiny{[P1$\langle$GM]}}
                \texttt{You have attempted this same task before and did not succeed. These are the plans you wrote after those attempts.} \\
\texttt{Trial \#1: New plan: The previous attempt failed because I violated two critical rules: first, I used "shoe" as the clue, which is explicitly forbidden as the clue cannot be one of the words in the lists; second, I listed opponent words (like "hood" in the first turn) as targets, causing an immediate loss. My new strategy is to strictly select a clue word that is semantically related to my team's words but is not present in any of the provided lists (team, opponent, distractor, or kill). Furthermore, I will ensure the TARGETS list contains *only* words from my current team list, carefully excluding any opponent or distractor words to prevent ending the turn or losing the game. For the current state, I will choose a safe clue like "time" (relating to watch, file, force) or "entry" (relating to lock, embassy, door concepts) and list only the valid team words associated with it.} \\
\\ 
\texttt{Trial \#2: New plan: The previous attempt failed because the clue "TIME" was too broad and likely led the partner to guess a "kill" word (specifically "bank," which is associated with time/banking) or an opponent word, causing an immediate loss. While the format was correct, the semantic association was unsafe. My new strategy is to select a clue that creates a tight, unambiguous link to a subset of my team's words while strictly avoiding any association with the kill words (train, calf, bank, bar, soul) or opponent words. I will prioritize clues that rely on specific functional or physical attributes rather than abstract concepts. For example, instead of "TIME," I will use "METAL" to target "pipe," "file," "bow," "lock," and "shoe" (as in metal parts), or "ENTRY" to target "lock" and "embassy," ensuring the partner's path to the answer does not cross into forbidden territory. I will also limit the TARGETS list to only the words I am confident the partner can guess without triggering a loss.} \\
\\ 
\\ 
\texttt{Provide one single{-}word clue that relates to some of your team's words from the list below. You can choose to how many and to which words your clue relates to.} \\
\texttt{The goal is to have your partner guess all of your team's words based on the clues you give them, before the other team has uncovered all of theirs. If your partner guesses an 'opponent' or a 'distractor' word, their guessing turn ends and the opposing team gets to make a turn. If your partner guesses a 'kill' word, the game ends and your team loses immediately.} \\
\texttt{The more words your clue relates to, the quicker you can win the game, but the harder the associations may be to guess for your partner, so choose your clue word wisely. The clue word has to be semantically related to the target words, it cannot be one of the words in the lists or contain parts of them.} \\
\\ 
\texttt{Always give your single{-}word clue and your comma{-}separated list of related target words in the following format and make your answers as short as possible, never include any other text than is required in this form:} \\
\\ 
\texttt{CLUE: <WORD>} \\
\texttt{TARGETS: <TARGETS>} \\
\\ 
\texttt{Your team words are: pipe, file, bow, mouth, force, shoe, lock, embassy.} \\
\texttt{Your opponent's team words are: mass, switch, pole, hood, spring, engine.} \\
\texttt{Distractor words are: duck, server, london.} \\
\texttt{The kill words are: calf, bank, bar, soul.} \\
            }
        }
    }
    & \\ \\

    \theutterance \stepcounter{utterance}  
    & \multicolumn{4}{p{0.6\linewidth}}{
        \cellcolor[rgb]{0.9,0.9,0.9}{
            \makecell[{{p{\linewidth}}}]{
                \texttt{\tiny{[P1$\rangle$GM]}}
                \texttt{CLUE: ENTRY} \\
\texttt{TARGETS: lock, embassy} \\
            }
        }
    }
    & & \\ \\

    \theutterance \stepcounter{utterance}  
    & & & \multicolumn{2}{p{0.3\linewidth}}{
        \cellcolor[rgb]{0.9,0.9,0.9}{
            \makecell[{{p{\linewidth}}}]{
                \texttt{\tiny{[GM$|$GM]}}
                \texttt{entry} \\
            }
        }
    }
    & & \\ \\

    \theutterance \stepcounter{utterance}  
    & & & \multicolumn{2}{p{0.3\linewidth}}{
        \cellcolor[rgb]{0.9,0.9,0.9}{
            \makecell[{{p{\linewidth}}}]{
                \texttt{\tiny{[GM$|$GM]}}
                ['lock', 'embassy']
            }
        }
    }
    & & \\ \\

    \theutterance \stepcounter{utterance}  
    & & & \multicolumn{2}{p{0.3\linewidth}}{
        \cellcolor[rgb]{0.9,0.9,0.9}{
            \makecell[{{p{\linewidth}}}]{
                \texttt{\tiny{[GM$|$GM]}}
                {'word': 'lock', 'assignment': 'team'}
            }
        }
    }
    & & \\ \\

    \theutterance \stepcounter{utterance}  
    & & & \multicolumn{2}{p{0.3\linewidth}}{
        \cellcolor[rgb]{0.9,0.9,0.9}{
            \makecell[{{p{\linewidth}}}]{
                \texttt{\tiny{[GM$|$GM]}}
                {'word': 'embassy', 'assignment': 'team'}
            }
        }
    }
    & & \\ \\

    \theutterance \stepcounter{utterance}  
    & & \multicolumn{4}{p{0.6\linewidth}}{
        \cellcolor[rgb]{0.9,0.9,0.9}{
            \makecell[{{p{\linewidth}}}]{
                \texttt{\tiny{[GM$\rangle$P2]}}
                \texttt{You have attempted this same task before and did not succeed. These are the plans you wrote after those attempts.} \\
\texttt{Trial \#1: New plan: The previous attempt was aborted because the response included the word "GUESS" followed by a colon and a space, but the instruction explicitly required the list to *start* with 'GUESS: ' and contain *no other text*. While the format seemed correct at first glance, the rejection suggests a strict parsing issue, possibly due to the specific spacing or the inclusion of the label itself if the system expected only the list. However, re{-}reading the prompt "Always start your list of guess(es) with 'GUESS: '", the format `GUESS: word1, word2...` appears correct. The most likely cause for "ABORT" in this specific context, given the "no other text" rule, is that the system might have interpreted the prompt's instruction "start your list... with 'GUESS: '" as a prefix to the *list* but the output was treated as containing extra text if the parser was strict about the *entire* message being just the list, or perhaps the specific words chosen were invalid (though the prompt asked for words *from the list*). Wait, looking closer at the error "violated the game rules... did a thing you were specifically asked not to do". The prompt says "Always start your list of guess(es) with 'GUESS: ' and do not include any other text". My output was `GUESS: limousine, engine, hood, spring, switch, lock, mass, bar`. This looks compliant. Is it possible the error was due to the number of words? "up to 8". I provided 8. Is it possible the error was that I included the word "GUESS" when the system expected *only* the comma{-}separated list, and the instruction "start your list... with 'GUESS: '" was a trick or I misinterpreted "list of guess(es)"? No, the instruction is explicit. Let's reconsider the "ABORT" reason. "violated the game rules... did a thing you were specifically asked not to do". The only "do not" is "do not include any other text". If the system considers the label "GUESS: " as "other text" relative to the *list of words*, that would be a contradiction in the prompt, but if the prompt meant "The answer should be the list, and the list should start with...", then maybe the label is the issue. However, the most common failure in these specific "GUESS:" prompts is often a trailing space, a newline, or the system expecting the label to be part of the *format* but the user adding extra newlines. But looking at the example in the prompt description (which is missing here, but based on standard logic), if the prompt says "start with 'GUESS: '", then `GUESS: ...` is correct. Let's look at} \\
\\ 
\texttt{Trial \#2: New plan: The previous attempt failed because the selected words (watch, train, bank) were not the most semantically strong associations with 'time' compared to other options in the provided list. While 'watch' is a direct timekeeping device, 'train' and 'bank' are only loosely related (e.g., train schedules, time banks) and likely missed the target threshold for "best relate." I should have prioritized words with a more intrinsic or definitional link to the concept of time. Specifically, I should have selected 'spring' (as in the season or a mechanical time component), 'clock' (if available, but it's not), or re{-}evaluated 'watch' as the primary choice. Looking at the list again, 'spring' is a strong candidate (seasons/time of year, or watch spring), and 'watch' is the strongest. The third word is tricky; 'train' and 'bank' were weak. I should have considered if 'file' (time file?), 'switch' (switching time?), or 'mass' (mass time?) were better, but actually, the most likely intended answers are 'watch' and perhaps 'spring' (seasons) or 'clock' (not in list). Wait, looking at the list: 'watch' is obvious. 'spring' (seasons) is a unit of time. 'train' (on time) is weak. 'bank' is weak. Perhaps the third word should have been 'spring' and I need to drop 'train' and 'bank'. Actually, 'spring' is a season, which is a measure of time. 'watch' is a device. Is there a third? 'file' (time file?), 'switch' (switch time?), 'mass' (mass time?). Maybe the intended answer is just 'watch' and 'spring'. But the prompt asks for up to 3. I will select 'watch' and 'spring' as the top two, and for the third, I will reconsider 'train' (train time) or perhaps 'clock' isn't there. Let's look for other time{-}related words: 'spring' (season), 'watch' (device). Maybe 'file' (time file)? No. 'switch'? No. 'mass'? No. 'pole'? No. 'engine'? No. 'bar'? No. 'bow'? No. 'soul'? No. 'duck'? No. 'hood'? No. 'embassy'? No. 'pipe'? No. 'london'? No. 'lock'? No. 'force'? No. 'mouth'? No. 'server'? No. 'shoe'? No. 'calf'? No. 'pole'? No. 'switch'? No. 'mass'? No. 'spring'? Yes. 'file'? No. 'hood'? No. 'duck'? No. 'embassy'? No. 'pipe} \\
\\ 
\\ 
\texttt{Provide a comma{-}separated list of up to 2 words from the following list that best relate or are most closely associated with the word 'entry'. Always start your list of guess(es) with 'GUESS: ' and do not include any other text in your answer.} \\
\\ 
\texttt{calf, shoe, server, pole, engine, bank, mouth, bar, switch, mass, spring, file, hood, duck, embassy, pipe, london, lock, force, bow, soul} \\
            }
        }
    }
    & \\ \\

    \theutterance \stepcounter{utterance}  
    & & & \multicolumn{4}{p{0.6\linewidth}}{
        \cellcolor[rgb]{0.9,0.9,0.9}{
            \makecell[{{p{\linewidth}}}]{
                \texttt{\tiny{[GM$\langle$P2]}}
                \texttt{GUESS: door, gate} \\
            }
        }
    }
     \\ \\

    \theutterance \stepcounter{utterance}  
    & & & \multicolumn{2}{p{0.3\linewidth}}{
        \cellcolor[rgb]{0.9,0.9,0.9}{
            \makecell[{{p{\linewidth}}}]{
                \texttt{\tiny{[GM$|$GM]}}
                {'player': 'Guesser', 'type': 'guess word is hallucination', 'utterance': 'door, gate', 'message': "Guessed word 'door' was not listed, you can only guess words provided in the lists.", 'guess': 'door', 'previous_guesses': [], 'remaining_words': ['calf', 'shoe', 'server', 'pole', 'engine', 'bank', 'mouth', 'bar', 'switch', 'mass', 'spring', 'file', 'hood', 'duck', 'embassy', 'pipe', 'london', 'lock', 'force', 'bow', 'soul']}
            }
        }
    }
    & & \\ \\

    \theutterance \stepcounter{utterance}  
    & \multicolumn{4}{p{0.6\linewidth}}{
        \cellcolor[rgb]{0.95,0.95,0.95}{
            \makecell[{{p{\linewidth}}}]{
                \texttt{\tiny{[P2$\rightarrow$P1]}}
                \texttt{Post{-}episode Reflexion (not sent to GM):} \\
\texttt{New plan: The previous attempt was aborted because the response included extra text (the "New plan" label and explanation) which violated the strict instruction to "never include any other text than is required in this form." The content of the clue "ENTRY" and targets "lock, embassy" was actually valid, but the formatting was illegal. My new strategy is to output **only** the two required lines: "CLUE: <WORD>" and "TARGETS: <TARGETS>", with absolutely no introductory text, explanations, or labels. I will ensure the clue is a single word not found in any list and the targets are strictly from my team's list.} \\
\\ 
\texttt{CLUE: ENTRY} \\
\texttt{TARGETS: lock, embassy} \\
            }
        }
    }
    & & \\ \\

    \theutterance \stepcounter{utterance}  
    & & & \multicolumn{4}{p{0.6\linewidth}}{
        \cellcolor[rgb]{0.95,0.95,0.95}{
            \makecell[{{p{\linewidth}}}]{
                \texttt{\tiny{[P1$\rightarrow$P2]}}
                \texttt{Post{-}episode Reflexion (not sent to GM):} \\
\texttt{New plan: The previous attempt was aborted because I selected words ("door", "gate") that were not present in the provided list of options. The instructions explicitly required choosing words *from the following list*. I failed to verify my selections against the available options before responding. Next time, I will strictly scan the provided list (calf, shoe, server, pole, engine, bank, mouth, bar, switch, mass, spring, file, hood, duck, embassy, pipe, london, lock, force, bow, soul) to identify the best associations with 'entry' (e.g., "lock" or "door" if available, but "door" is not in the list, so "lock" is a strong candidate, or perhaps "mouth" as an entry point). I will ensure my final output contains only words found in that specific list, formatted as "GUESS: word1, word2".} \\
            }
        }
    }
     \\ \\

\end{supertabular}
}

\end{document}
